\chapter{Vector Calculus: Green, Stokes, and Divergence}

\providecommand{\oiint}{\oint\kern-0.62em\oint}

In the previous chapter on multiple integrals, we learned how to compute double
integrals $\iint_D f\,\dd A$ over flat regions and triple integrals $\iiint_V f\,\dd V$
over solid shapes. In every case, we chopped a region into a fine grid of cells and
summed a scalar quantity---like density, height, or charge---throughout the entire
interior. Every single point inside contributed its share to the total.

What happens if the quantity filling the interior is not an arbitrary scattering of
numbers, but the rate of change of an underlying field? Every theorem in this chapter
expresses the same answer: accumulating a derivative over the interior of a region
equals evaluating the field itself along the boundary.

Back in Chapter~21, the Fundamental Theorem of Calculus revealed a remarkable
simplification in one dimension: to add up a rate of change over an interval, you
do not actually have to add up anything in the interior at all -- you just look at
the two endpoints:
\[
  \int_a^b F'(x)\,\dd x \;=\; F(b) - F(a).
\]
The whole interior of $[a,b]$ -- the infinitely many places where $F$ is wiggling
and changing -- collapses into two numbers at the boundary. The accumulated change
inside is completely determined by the values at the edge.

That principle -- \emph{the boundary records what happens inside} -- is the central
theme of vector calculus. The great theorems of this chapter (Green's theorem, the
divergence theorem, Stokes' theorem) all say the same thing in higher dimensions:
to learn what is happening throughout a multi-dimensional region, you only have to
walk around its \emph{boundary}. An integral over a solid blob equals an integral
over the skin that wraps it. An integral over a flat patch equals an integral around
the loop that fences it. In each case, integrating a derivative over a $d$-dimensional region
equals evaluating the field along its $(d{-}1)$-dimensional boundary. The Fundamental Theorem of Calculus
is the one-dimensional version, where the boundary of an interval
consists of its two endpoints.

To formulate these theorems we use \emph{vector fields}. Up to this point our
functions have assigned numbers to points. A vector field assigns a
\emph{vector} to each point: a velocity at every point of a flowing fluid, a
gravitational or electrostatic force, or the flow of heat. Once we can
integrate vector fields along curves and across surfaces, the integral theorems
take a clean, unified form.

\begin{focusbox}
One idea, four faces:
\[
\begin{aligned}
\text{FTC:} &\quad \int_a^b F' = F(b)-F(a) &&\text{(1D: endpoints)}\\
\text{Gradient theorem:} &\quad \int_C \grad f\cdot \dd\vec r = f(\text{end})-f(\text{start}) &&\text{(curve: endpoints)}\\
\text{Green / Stokes:} &\quad \iint_S (\text{curl}) = \oint_{\partial S}(\text{field}) &&\text{(patch: boundary loop)}\\
\text{Divergence:} &\quad \iiint_V (\text{div}) = \oiint_{\partial V}(\text{flux}) &&\text{(solid: boundary skin)}
\end{aligned}
\]
Read every line as: \emph{accumulate a derivative inside} $=$ \emph{read the
field on the boundary}.
\end{focusbox}

\chapterline

\section{Vector fields: assigning vectors to points}
\label{ch35:fields}

A \emph{scalar field} assigns a number to each point: temperature in a room,
altitude on a map, pressure in the ocean. We have been studying these for two
chapters under the name $f(x,y)$ or $f(x,y,z)$. A \emph{vector field} assigns an
\emph{arrow} -- a vector -- to each point instead.

\begin{definition}[Vector field]
\label{ch35:def-field}
A vector field on a region of the plane is a function
\[
  \vec F(x,y) = \big(P(x,y),\, Q(x,y)\big) = P(x,y)\,\vec i + Q(x,y)\,\vec j,
\]
that assigns to each point $(x,y)$ the vector $\vec F(x,y)$. In three dimensions,
$\vec F(x,y,z) = (P,Q,R)$, with three component functions. We draw it by planting
the arrow $\vec F(x,y)$ with its \emph{tail} at the point $(x,y)$.
\end{definition}

The geometric picture is direct. Below is the field $\vec F(x,y) = (-y, x)$. At each
point we draw the arrow $(-y,x)$. Notice it always points ninety degrees
counterclockwise from the position vector $(x,y)$ -- so the whole field swirls
around the origin like water circling a drain, or a spinning turntable.

\begin{figure}[H]
\centering
\begin{tikzpicture}[scale=0.9,>={Stealth[length=4pt]}]
% Left Panel: Rotational Swirl
\begin{scope}[shift={(0,0)}]
  \draw[fill=blue!3,draw=blue!15,rounded corners=4pt] (-2.7,-2.7) rectangle (2.7,2.7);
  \node[font=\small\bfseries,anchor=north] at (0,2.55) {Pure Circulation: $\vec F=(-y,x)$};
  \draw[->,gray!50] (-2.3,0)--(2.3,0) node[right]{\scriptsize $x$};
  \draw[->,gray!50] (0,-2.3)--(0,2.3) node[above]{\scriptsize $y$};
  % swirl guides
  \draw[curveblue!25,dashed,thick] (0,0) circle (1.0);
  \draw[curveblue!25,dashed,thick] (0,0) circle (1.8);
  \foreach \x in {-1.8,-0.9,0.9,1.8}{
    \foreach \y in {-1.8,-0.9,0.9,1.8}{
      \draw[->,curveblue,thick] (\x,\y) -- ++({-\y*0.22},{\x*0.22});
    }
  }
  \node[curveblue!80!black,font=\scriptsize,fill=white,inner sep=1.5pt,rounded corners=2pt] at (0,-1.4) {$\operatorname{curl} = 2,\;\operatorname{div} = 0$};
\end{scope}

% Right Panel: Radial Source
\begin{scope}[shift={(6.2,0)}]
  \draw[fill=red!3,draw=red!15,rounded corners=4pt] (-2.7,-2.7) rectangle (2.7,2.7);
  \node[font=\small\bfseries,anchor=north] at (0,2.55) {Pure Outflow: $\vec F=(x,y)$};
  \draw[->,gray!50] (-2.3,0)--(2.3,0) node[right]{\scriptsize $x$};
  \draw[->,gray!50] (0,-2.3)--(0,2.3) node[above]{\scriptsize $y$};
  % expansion rings
  \draw[tangentred!25,dashed,thick] (0,0) circle (1.0);
  \draw[tangentred!25,dashed,thick] (0,0) circle (1.8);
  \foreach \x in {-1.8,-0.9,0.9,1.8}{
    \foreach \y in {-1.8,-0.9,0.9,1.8}{
      \draw[->,tangentred,thick] (\x,\y) -- ++({\x*0.22},{\y*0.22});
    }
  }
  \node[tangentred!80!black,font=\scriptsize,fill=white,inner sep=1.5pt,rounded corners=2pt] at (0,-1.4) {$\operatorname{div} = 2,\;\operatorname{curl} = 0$};
\end{scope}
\end{tikzpicture}
\caption{The two fundamental motions of a vector field. Left: The rotational field
$\vec F(x,y)=(-y,x)$ circulates around the origin with constant scalar curl $2$ and
zero divergence. Right: The source field $\vec F(x,y)=(x,y)$ radiates outward from the
origin with constant divergence $2$ and zero curl.}
\label{ch35:fig-swirl}
\end{figure}

\begin{intuitionbox}
Two motions live inside every vector field, and the rest of the chapter is about
measuring them. \textbf{Circulation}: does the field tend to \emph{spin} you
around a loop? (The swirl in Figure~\ref{ch35:fig-swirl}, left, does.) \textbf{Flux}:
does the field tend to \emph{carry stuff out through} a boundary? (The source in
Figure~\ref{ch35:fig-swirl}, right, does.) Curl will measure local spin; divergence
will measure local spreading. The integral theorems make this relationship quantitative.
\end{intuitionbox}

\begin{recoverybox}
\textbf{Pause and think}: Look closely at the arrows in Figure~\ref{ch35:fig-swirl}.
If you drop a tiny floating twig into the swirl $\vec F = (-y, x)$, it spins around
its own center of mass as it circles the origin. But if you drop the same twig into
the radial outflow $\vec F = (x, y)$, it is pushed straight outward along a ray,
never turning. Spin versus spread: this chapter is entirely about how these two simple
local tendencies add up across curves, surfaces, and solids.
\end{recoverybox}

\begin{applicationbox}
Vector fields model physical quantities distributed through space. The wind on a
weather map is $\vec F(x,y)=$ (velocity of air). The gravitational pull near a
planet is a field of arrows pointing inward. The electric field around a charge,
the magnetic field around a wire, the flow of heat through a wall, the current in
the ocean -- each is a vector field, and each obeys laws written in the language
of this chapter. Maxwell's four equations of electromagnetism are statements
about the divergence and curl of electric and magnetic fields.
\end{applicationbox}

\section{Line integrals: adding a field along a path}
\label{ch35:lineint}

\subsection{Work, and why we multiply force by displacement}

Drag a box across a floor. The work you do is force times distance -- but only the
\emph{part of the force that lies along the motion} counts. Push straight down on
a box and slide it sideways and you have done no work on it at all; your downward
force is perpendicular to the sideways motion. The right measure of work for a
constant force $\vec F$ over a straight displacement $\vec d$ is the dot product
\[
  W = \vec F \cdot \vec d = \abs{\vec F}\,\abs{\vec d}\cos\theta,
\]
which keeps exactly the component of $\vec F$ pointing along $\vec d$ and throws
away the perpendicular part.

Now let the force vary from point to point -- a vector field $\vec F$ -- and let
the path curve. We do what we always do in calculus: chop the path into tiny
nearly-straight pieces, on each piece treat the force as roughly constant, compute
the little bit of work $\vec F\cdot\Delta\vec r$, add them all up, and take the
limit as the pieces shrink. That limit is the \emph{line integral}.

\begin{definition}[Line integral of a vector field]
\label{ch35:def-lineint}
Let $C$ be a curve traced by $\vec r(t)=(x(t),y(t))$ for $t\in[a,b]$, and let
$\vec F$ be a vector field along it. The line integral of $\vec F$ along $C$ is
\[
  \int_C \vec F\cdot \dd\vec r
  \;=\; \int_a^b \vec F\big(\vec r(t)\big)\cdot \vec r\,'(t)\,\dd t.
\]
Writing $\vec F=(P,Q)$ and $\dd\vec r = (\dd x,\dd y)$, this is also written
$\displaystyle \int_C P\,\dd x + Q\,\dd y$.
\end{definition}

\begin{figure}[H]
\centering
\begin{tikzpicture}[scale=1.0,>={Stealth[length=5pt]}]
  % A smooth curve C
  \draw[curveblue, very thick] (0,0.5) .. controls (2,2.3) and (4,0.3) .. (7.5,2.4)
    node[pos=0.96, below right, font=\small\bfseries] {$C$};
  % Point of interest
  \coordinate (P) at (2.4, 1.48);
  \fill[black] (P) circle (2pt) node[below=4pt, font=\scriptsize] {$\vec r(t)$};
  % Tangent displacement vector dr
  \draw[->, sourceblue, ultra thick] (P) -- ++(2.2, 0.55) coordinate (dr)
    node[right, font=\scriptsize\bfseries] {$\dd\vec r = \vec r\,'(t)\,\dd t$};
  % Force field vector F
  \draw[->, tangentred, ultra thick] (P) -- ++(1.4, 1.7) coordinate (F)
    node[above right, font=\scriptsize\bfseries] {$\vec F(\vec r(t))$};
  % Angle arc
  \draw[gray!60, thick] (P) ++(0.7, 0.175) arc (14:50:0.7);
  \node[font=\scriptsize] at (3.2, 1.85) {$\theta$};
  % Projection dashed line
  \coordinate (proj) at (3.75, 1.82);
  \draw[dashed, gray!70, thick] (F) -- (proj);
  % Right angle indicator
  \draw[gray!70] (proj) ++(-0.06, 0.22) -- ++(0.22, 0.05) -- ++(0.06, -0.22);
  % Projected component brace
  \draw[decorate, decoration={brace, amplitude=4pt, mirror, raise=4pt}, quotientorange, thick]
    (P) -- (proj) node[midway, below=8pt, font=\scriptsize\bfseries] {$|\vec F|\cos\theta$};
  % Explanatory text card
  \node[fill=blue!3, draw=blue!20, rounded corners=3pt, font=\scriptsize, align=left, inner sep=5pt]
    at (1.1, 3.1) {\textbf{Infinitesimal Work}:\\$\dd W = \vec F \cdot \dd\vec r = |\vec F||\dd\vec r|\cos\theta$\\$\hphantom{\dd W} = P\,\dd x + Q\,\dd y$};
\end{tikzpicture}
\caption{The infinitesimal work element $\vec F\cdot\dd\vec r$ along a curve $C$. At each
point $\vec r(t)$, the path advances along the tangent displacement $\dd\vec r = \vec r\,'(t)\,\dd t$.
The dot product projects $\vec F$ onto the motion: only the component $|\vec F|\cos\theta$ parallel
to the curve does work. The perpendicular component produces zero motion along $C$.}
\label{ch35:fig-lineint-tangent}
\end{figure}

\begin{whybox}
Why does $\dd\vec r$ become $\vec r\,'(t)\,\dd t$? Because $\vec r\,'(t)$ is the
velocity vector -- it points along the curve and its length is the speed. Over a
tiny time $\dd t$ the path advances by the little displacement vector
$\dd\vec r = \vec r\,'(t)\,\dd t$. Dotting the force against this displacement
gives the infinitesimal work $\vec F\cdot\dd\vec r$, and integrating in $t$ adds
them along the curve. The parameter $t$ traces the path; the
value of the integral does not depend on the parametrization speed.
\end{whybox}

\begin{walkbox}
Why doesn't the speed of traversal change the answer? If you re-parametrize and
move twice as fast, then $\vec r\,'$ doubles but the time interval $\dd t$ halves,
and the product $\vec r\,'(t)\,\dd t = \dd\vec r$ is unchanged -- it is a geometric
displacement, independent of the rate of traversal. (Reversing direction, however, flips every
$\dd\vec r$, so it flips the sign of the integral. Orientation matters; speed does
not.)
\end{walkbox}

\subsection{A worked line integral, step by step}

\begin{example}[Work along a parabola]
\label{ch35:ex-parab}
Compute $\displaystyle\int_C \vec F\cdot\dd\vec r$ for $\vec F(x,y)=(y,\,x^2)$
along the parabola $y=x^2$ from $(0,0)$ to $(2,4)$.

\textbf{Step 1 -- parametrize.} Let $x=t$, $y=t^2$, so $\vec r(t)=(t,t^2)$ for
$t\in[0,2]$. Then $\vec r\,'(t)=(1,2t)$.

\textbf{Step 2 -- evaluate the field on the path.}
$\vec F(\vec r(t)) = (y,\,x^2)\big|_{x=t,\,y=t^2} = (t^2,\, t^2)$.

\textbf{Step 3 -- dot.}
$\vec F(\vec r(t))\cdot\vec r\,'(t) = (t^2)(1)+(t^2)(2t) = t^2 + 2t^3.$

\textbf{Step 4 -- integrate.}
\[
  \int_0^2 (t^2 + 2t^3)\,\dd t
  = \left[\frac{t^3}{3} + \frac{t^4}{2}\right]_0^2
  = \frac{8}{3} + 8 = \frac{32}{3}.
\]
The total work is $\tfrac{32}{3}$. The recipe never changes: parametrize, plug in,
dot, integrate a single ordinary integral in $t$.
\end{example}

\begin{figure}[H]
\centering
\begin{tikzpicture}[scale=1.0,>={Stealth[length=4pt]}]
\begin{axis}[
  width=8.4cm, height=6.4cm,
  axis lines=middle, xlabel={\scriptsize $x$}, ylabel={\scriptsize $y$},
  xmin=-0.4, xmax=2.6, ymin=-0.4, ymax=4.6,
  xtick={1,2}, ytick={1,2,3,4}, tick label style={font=\scriptsize},
  clip=false,
]
% the path y = x^2
\addplot[curveblue, very thick, domain=0:2, samples=80] {x^2};
\node[curveblue, anchor=west] at (axis cs:1.55,2.5) {\scriptsize $C:\ y=x^2$};
% field arrows F=(y,x^2)=(t^2,t^2) at sample points (t,t^2) along the path
\addplot[tangentred, ->, quiver={u=\thisrow{u}, v=\thisrow{v}}, only marks, mark=none]
  table[row sep=\\]{
    x    y    u     v    \\
    0.5  0.25 0.075 0.075 \\
    1.0  1.0  0.30  0.30  \\
    1.5  2.25 0.675 0.675 \\
    2.0  4.0  1.20  1.20  \\
  };
% endpoints
\addplot[only marks, mark=*, mark size=1.4pt, black] coordinates {(0,0) (2,4)};
\node[anchor=north west] at (axis cs:0.05,-0.05) {\scriptsize $(0,0)$};
\node[anchor=west] at (axis cs:2.05,4) {\scriptsize $(2,4)$};
\end{axis}
\end{tikzpicture}
\caption{The work integral of Example~\ref{ch35:ex-parab}: we walk the parabola
$C$ (blue) from $(0,0)$ to $(2,4)$. The line integral adds up, at each point of the
path, the component of the field $\vec F=(y,x^2)$ that lies along the blue curve.
The along-path total works out to $\tfrac{32}{3}$.}
\label{ch35:fig-work}
\end{figure}

\begin{trapbox}
A line integral $\int_C \vec F\cdot\dd\vec r$ is \emph{not} the same as the
arc-length integral $\int_C f\,\dd s$ you may have seen (mass of a wire). The dot
product version measures the field's component \emph{along} the curve and is
sensitive to direction; reversing the path negates it. The arc-length version uses
$\dd s = \abs{\vec r\,'(t)}\,\dd t$ (a length, always positive) and does not care
about direction. They answer different questions. This chapter is almost entirely
about the dot-product kind.
\end{trapbox}

\section{Conservative fields: when the path does not matter}
\label{ch35:conservative}

In Example~\ref{ch35:ex-parab} the answer depended on the curve we chose. For some
special fields it does \emph{not}: the work depends only on where you start and
where you stop, never on the route between. These are the \emph{conservative}
fields, and they are exactly the fields that are gradients.

\begin{theorem}[Fundamental Theorem for line integrals / Gradient theorem]
\label{ch35:thm-gradient}
If $\vec F = \grad f$ for some scalar function $f$ (a \emph{potential}), then for
any curve $C$ from point $A$ to point $B$,
\[
  \int_C \grad f\cdot \dd\vec r = f(B) - f(A).
\]
The integral depends only on the endpoints. In particular, around any closed loop
($A=B$), $\displaystyle\oint_C \grad f\cdot\dd\vec r = 0$.
\end{theorem}

\begin{whybox}
This is the Fundamental Theorem of Calculus applied along a curve. Parametrize $C$ by
$\vec r(t)$, $t\in[a,b]$, and watch the chain rule do the entire job:
\[
  \int_C \grad f\cdot\dd\vec r
  = \int_a^b \grad f(\vec r(t))\cdot \vec r\,'(t)\,\dd t
  = \int_a^b \frac{\dd}{\dd t}\Big[f(\vec r(t))\Big]\,\dd t
  = f(\vec r(b)) - f(\vec r(a)).
\]
The middle equality is the multivariable chain rule (Chapter~33):
$\grad f\cdot\vec r\,' = \dv{}{t}f(\vec r(t))$. So the integrand was an exact
derivative, and the ordinary FTC finishes it. Here, the value of the integral
depends only on the potential at the two endpoints.
\end{whybox}

\begin{intuitionbox}
A potential $f$ is an \emph{altitude}. The field $\grad f$ points straight uphill,
and its work is the height you gain. Walk any trail from the trailhead to the
summit -- switchbacks, scenic detours, whatever -- and the elevation gained is the
same, because it is just (summit altitude) minus (trailhead altitude). Walk a loop
back to where you started and you gain zero net altitude. That is the whole content
of conservativeness. A field is conservative precisely when it has an altitude
function.
\end{intuitionbox}

\begin{figure}[H]
\centering
\begin{tikzpicture}[scale=1.0,>={Stealth[length=4pt]}]
\begin{axis}[
  width=8.6cm, height=6.6cm,
  axis lines=middle, xlabel={\scriptsize $x$}, ylabel={\scriptsize $y$},
  xmin=-2.3, xmax=2.3, ymin=-2.3, ymax=2.3,
  xtick=\empty, ytick=\empty,
]
% level curves of the potential f = x^2 + y^2 (circles)
\addplot[gray!60, samples=90, domain=0:360] ({0.7*cos(x)},{0.7*sin(x)});
\addplot[gray!60, samples=90, domain=0:360] ({1.3*cos(x)},{1.3*sin(x)});
\addplot[gray!60, samples=90, domain=0:360] ({1.9*cos(x)},{1.9*sin(x)});
% gradient arrows F = grad f = (2x,2y), scaled, pointing outward
\addplot[targetgreen, ->, quiver={u=\thisrow{u}, v=\thisrow{v}}, only marks, mark=none]
  table[row sep=\\]{
    x     y     u      v     \\
    -1.5  -1.5  -0.42  -0.42 \\
    -0.8  -1.5  -0.224 -0.42 \\
     0.8  -1.5   0.224 -0.42 \\
     1.5  -1.5   0.42  -0.42 \\
    -1.5  -0.8  -0.42  -0.224 \\
    -0.8  -0.8  -0.224 -0.224 \\
     0.8  -0.8   0.224 -0.224 \\
     1.5  -0.8   0.42  -0.224 \\
    -1.5   0.8  -0.42   0.224 \\
    -0.8   0.8  -0.224  0.224 \\
     0.8   0.8   0.224  0.224 \\
     1.5   0.8   0.42   0.224 \\
    -1.5   1.5  -0.42   0.42 \\
    -0.8   1.5  -0.224  0.42 \\
     0.8   1.5   0.224  0.42 \\
     1.5   1.5   0.42   0.42 \\
  };
\node[gray!50!black, anchor=west] at (axis cs:0.95,1.7) {\scriptsize level sets of $f$};
\node[targetgreen!60!black, anchor=west] at (axis cs:1.0,-2.0) {\scriptsize $\grad f$};
\end{axis}
\end{tikzpicture}
\caption{A conservative field is a gradient. Here $f(x,y)=x^2+y^2$ is the
``altitude.'' Its level sets (gray circles) are the contour lines, and the field
$\grad f=(2x,2y)$ (green arrows) crosses them at right angles, always pointing
uphill. Work done along any path is just the change in altitude between its
endpoints; around a closed loop it is zero.}
\label{ch35:fig-gradient}
\end{figure}

\subsection{How to test for, and find, a potential}

In the plane, write $\vec F=(P,Q)$. If $\vec F=\grad f=(f_x,f_y)$ then
$P=f_x$ and $Q=f_y$. By Clairaut's theorem (mixed partials commute, Chapter~32),
\[
  \pdv{P}{y} = f_{xy} = f_{yx} = \pdv{Q}{x}.
\]
So a \emph{necessary} condition is $P_y = Q_x$. On a region with no holes (a
``simply connected'' region -- one where every loop can be shrunk to a point) this
condition is also \emph{sufficient}: it guarantees a potential exists.

\begin{example}[Finding a potential]
\label{ch35:ex-potential}
Is $\vec F=(2xy + 3,\; x^2 - 1)$ conservative? Here $P=2xy+3$, $Q=x^2-1$.

\textbf{Test.} $P_y = 2x$ and $Q_x = 2x$. They agree, and the plane has no holes,
so a potential exists.

\textbf{Build it.} We need $f$ with $f_x = 2xy+3$. Integrate in $x$ (treating $y$
as constant):
\[
  f(x,y) = x^2 y + 3x + g(y),
\]
where the ``constant'' of integration is allowed to depend on $y$. Now demand
$f_y = Q$:
\[
  f_y = x^2 + g'(y) \overset{!}{=} x^2 - 1 \;\Rightarrow\; g'(y) = -1 \;\Rightarrow\; g(y) = -y + C.
\]
So $f(x,y) = x^2y + 3x - y + C$. As a check, $\grad f = (2xy+3,\; x^2-1)=\vec F$.
Now any line integral of $\vec F$ is trivial: from $(0,0)$ to $(1,2)$ it is
$f(1,2)-f(0,0) = (2 + 3 - 2) - 0 = 3$, no matter the path.
\end{example}

\begin{trapbox}
The hole matters. The famous swirl
$\vec F = \big(\frac{-y}{x^2+y^2}, \frac{x}{x^2+y^2}\big)$ satisfies $P_y=Q_x$
\emph{everywhere it is defined}, yet $\oint \vec F\cdot\dd\vec r = 2\pi \ne 0$
around the unit circle. It is \emph{not} conservative. The catch: it is undefined
at the origin, so its domain has a hole, and the loop around that hole cannot be
shrunk to a point. The test $P_y=Q_x$ only promises a potential on a
\emph{simply connected} domain. Here topology enters multivariable calculus.
\end{trapbox}

\begin{connectionbox}
The word ``conservative'' comes from physics, not bookkeeping. If $\vec F$ is a
force with potential energy $U$ (so $\vec F=-\grad U$), then the work done equals
the drop in $U$, and the total mechanical energy (kinetic plus potential) stays
strictly constant along any trajectory -- energy is \emph{conserved}. Gravity
and electrostatics are conservative; friction is not (it bleeds energy away as
heat, and its closed-loop integral is non-zero, which is why perpetual-motion
machines cannot work). The nineteenth-century drive to understand which physical
forces conserve mechanical energy, spearheaded by Hermann von Helmholtz and Lord
Kelvin, is precisely what led mathematicians to clarify when a vector field is the
gradient of an altitude function.
\end{connectionbox}

\section{Green's theorem: interior curl equals boundary circulation}
\label{ch35:green}

The first of the three great theorems is also the one easiest to picture.
Green's theorem relates a line integral \emph{around} a closed loop $C$ to a double
integral \emph{over} the flat region $D$ that the loop encloses.

\begin{theorem}[Green's theorem, circulation form]
\label{ch35:thm-green}
Let $D$ be a region in the plane bounded by a simple closed curve $C$, traversed
counterclockwise. If $\vec F=(P,Q)$ has continuous partials on $D$, then
\[
  \oint_C P\,\dd x + Q\,\dd y \;=\; \iint_D \left(\pdv{Q}{x} - \pdv{P}{y}\right)\dd A.
\]
\end{theorem}

The quantity $\pdv{Q}{x}-\pdv{P}{y}$ is the \emph{microscopic circulation} -- the
``curl'' of the field, which we meet properly in the next section. Green's theorem
says: the total swirl around the boundary equals the sum of all the little swirls
inside.

\begin{happeningbox}
The underlying mechanism is a telescoping cancellation. Tile the region $D$ with tiny
rectangles. For each tiny rectangle, run the line integral around its own little
boundary, counterclockwise. Now add up all these little loop integrals. Look at
any internal edge shared by two adjacent tiles: one tile traverses it
left-to-right, its neighbor traverses the \emph{same} edge right-to-left. The two
contributions are equal and opposite -- they \textbf{cancel}. Every interior edge
cancels against its neighbor. The only edges that survive this
cancellation are those on the outer boundary of $D$, where there is no adjacent tile
to cancel with. So
\[
  \sum_{\text{tiles}} (\text{tiny loop integral}) = (\text{big boundary integral}).
\]
Meanwhile each tiny loop integral, by a short computation, equals
$\big(Q_x - P_y\big)\times(\text{area of tile})$. Summing those is the double
integral. The interior cancellation is \emph{exactly} the same telescoping that
made the Fundamental Theorem of Calculus work, where the inner endpoints of
adjacent subintervals cancelled and only $a$ and $b$ survived.
\end{happeningbox}

\begin{figure}[H]
\centering
\begin{tikzpicture}[scale=0.92,>={Stealth[length=4pt]}]
  % Panel 1: Macroscopic Loop
  \begin{scope}[shift={(0,0)}]
    \draw[fill=blue!3,draw=blue!15,rounded corners=4pt] (-2.2,-2.2) rectangle (2.2,2.2);
    \node[font=\small\bfseries,anchor=north] at (0,2.05) {(a) Macro boundary $\partial D$};
    \filldraw[fill=blue!10,draw=curveblue,very thick]
      (-1.3,-1.3) rectangle (1.3,1.3);
    \draw[->,curveblue,ultra thick] (-1.3,-1.3) -- (0.2,-1.3);
    \draw[->,curveblue,ultra thick] (1.3,-1.3) -- (1.3,0.2);
    \draw[->,curveblue,ultra thick] (1.3,1.3) -- (-0.2,1.3);
    \draw[->,curveblue,ultra thick] (-1.3,1.3) -- (-1.3,-0.2);
    \node[curveblue!80!black,font=\small\bfseries] at (0,0) {$D$};
    \node[curveblue,font=\scriptsize] at (1.5,0.7) {$C=\partial D$};
    \node[font=\scriptsize,align=center] at (0,-1.75) {$\oint_C \vec F\cdot\dd\vec r$};
  \end{scope}

  % Panel 2: Microscopic Circulation Grid
  \begin{scope}[shift={(5.0,0)}]
    \draw[fill=blue!3,draw=blue!15,rounded corners=4pt] (-2.2,-2.2) rectangle (2.2,2.2);
    \node[font=\small\bfseries,anchor=north] at (0,2.05) {(b) Sum of micro tiles};
    % 3x3 grid
    \foreach \gx in {-1.2,-0.4,0.4}{
      \foreach \gy in {-1.2,-0.4,0.4}{
        \draw[fill=white,draw=gray!40] (\gx,\gy) rectangle ++(0.8,0.8);
        % micro ccw loop
        \draw[->,curveblue!70,thick] (\gx+0.15,\gy+0.15) -- (\gx+0.65,\gy+0.15);
        \draw[->,curveblue!70,thick] (\gx+0.65,\gy+0.15) -- (\gx+0.65,\gy+0.65);
        \draw[->,curveblue!70,thick] (\gx+0.65,\gy+0.65) -- (\gx+0.15,\gy+0.65);
        \draw[->,curveblue!70,thick] (\gx+0.15,\gy+0.65) -- (\gx+0.15,\gy+0.15);
      }
    }
    \node[font=\scriptsize,align=center] at (0,-1.75) {$\sum_{\text{tiles}} \oint_{\partial T_i} \vec F\cdot\dd\vec r$};
  \end{scope}

  % Panel 3: Internal Cancellation
  \begin{scope}[shift={(10.0,0)}]
    \draw[fill=blue!3,draw=blue!15,rounded corners=4pt] (-2.2,-2.2) rectangle (2.2,2.2);
    \node[font=\small\bfseries,anchor=north] at (0,2.05) {(c) Edge cancellation};
    % Two adjacent tiles
    \draw[fill=white,draw=gray!40] (-1.3,-0.6) rectangle (-0.05,0.6);
    \draw[fill=white,draw=gray!40] (0.05,-0.6) rectangle (1.3,0.6);
    % Shared internal edge opposing arrows
    \draw[->,tangentred,very thick] (-0.05,-0.4) -- (-0.05,0.4);
    \draw[->,tangentred,very thick] (0.05,0.4) -- (0.05,-0.4);
    \node[tangentred!90!black,font=\scriptsize\bfseries,align=center] at (0,0.9) {opposing $\implies$\\cancels!};
    % Outer surviving edges
    \draw[->,curveblue,very thick] (-1.3,-0.6) -- (-0.05,-0.6);
    \draw[->,curveblue,very thick] (0.05,-0.6) -- (1.3,-0.6);
    \draw[->,curveblue,very thick] (1.3,-0.6) -- (1.3,0.6);
    \draw[->,curveblue,very thick] (1.3,0.6) -- (0.05,0.6);
    \draw[->,curveblue,very thick] (-0.05,0.6) -- (-1.3,0.6);
    \draw[->,curveblue,very thick] (-1.3,0.6) -- (-1.3,-0.6);
    \node[font=\scriptsize,align=center] at (0,-1.75) {Only boundary $\partial D$\\survives};
  \end{scope}
\end{tikzpicture}
\caption{Green's theorem via microscopic edge cancellation. (a) The macro circulation
$\oint_{\partial D} \vec F\cdot\dd\vec r$ around the outer boundary loop. (b) The region
$D$ tiled into tiny sub-rectangles, each evaluated counterclockwise. (c) On every internal
boundary shared by two adjacent cells, one cell traverses the segment forward while the
other traverses it backward; the two opposing contributions cancel out identically.
Every interior edge cancels, leaving only the outer perimeter $\partial D$ to account for
the sum of all microscopic curls.}
\label{ch35:fig-green-cancel}
\end{figure}

\subsection{Computing area with Green's theorem}

As an immediate application, choose $\vec F$ so that $Q_x - P_y = 1$. Then the double
integral is just the \emph{area} of $D$, and Green's theorem computes area from a
trip around the boundary. Taking $\vec F=(-y/2,\,x/2)$ gives
$Q_x-P_y=\tfrac12+\tfrac12=1$, so
\[
  \text{Area}(D) = \frac{1}{2}\oint_C x\,\dd y - y\,\dd x.
\]
This is the principle behind the \emph{planimeter}, a mechanical instrument that
measures the area of any shape by tracing its outline -- once around the boundary
and a dial reads off the enclosed area.

\begin{example}[Area of an ellipse by its boundary]
\label{ch35:ex-ellipse}
Let $C$ be the ellipse $x=a\cos t$, $y=b\sin t$, $t\in[0,2\pi]$. Then
$\dd x = -a\sin t\,\dd t$, $\dd y = b\cos t\,\dd t$, and
\[
  x\,\dd y - y\,\dd x
  = (a\cos t)(b\cos t)\,\dd t - (b\sin t)(-a\sin t)\,\dd t
  = ab(\cos^2 t + \sin^2 t)\,\dd t = ab\,\dd t.
\]
Therefore
\[
  \text{Area} = \frac12\int_0^{2\pi} ab\,\dd t = \frac12\,(ab)(2\pi) = \pi ab.
\]
This yields the familiar $\pi a b$ using only the boundary curve.
\end{example}

\begin{example}[Verifying Green's theorem directly]
\label{ch35:ex-green-verify}
Let $\vec F=(-y,\,x)$ (the swirl) and let $D$ be the unit disk, $C$ its boundary
circle counterclockwise.

\textbf{Right side (double integral).} Here $P=-y$, $Q=x$, so
$Q_x - P_y = 1 - (-1) = 2$. Thus
$\iint_D 2\,\dd A = 2\cdot(\text{area of disk}) = 2\pi.$

\textbf{Left side (line integral).} Parametrize $x=\cos t$, $y=\sin t$:
\[
  \oint_C -y\,\dd x + x\,\dd y
  = \int_0^{2\pi}\big[(-\sin t)(-\sin t) + (\cos t)(\cos t)\big]\dd t
  = \int_0^{2\pi} 1\,\dd t = 2\pi.
\]
Both sides give $2\pi$. The constant microscopic swirl of $2$ at every point,
summed over the disk, equals the circulation you feel going once around the rim.
\end{example}

\section{Curl and divergence: spin and spread, made local}
\label{ch35:curl-div}

Green's theorem has already handed us two local quantities worth naming. They are
the derivatives the great theorems integrate, so it pays to know what each one
measures physically.

\subsection{Divergence: how much a point is a source}

\begin{definition}[Divergence]
\label{ch35:def-div}
For $\vec F=(P,Q,R)$, the divergence is the scalar
\[
  \operatorname{div}\vec F = \grad\cdot\vec F = \pdv{P}{x} + \pdv{Q}{y} + \pdv{R}{z}.
\]
(In the plane, drop the last term.)
\end{definition}

\begin{intuitionbox}
Divergence measures \emph{spreading}. Imagine $\vec F$ is the velocity of a fluid.
Surround a point with a tiny balloon. If more fluid flows \emph{out} of the balloon
than in, the point is a \emph{source} and $\operatorname{div}\vec F>0$. If fluid piles
\emph{in}, it is a \emph{sink} and $\operatorname{div}\vec F<0$. If as much leaves as
enters, $\operatorname{div}\vec F=0$ and the fluid is \emph{incompressible} there.
For the source field $\vec F=(x,y)$ of Figure~\ref{ch35:fig-swirl} (right),
$\operatorname{div}\vec F = 1+1 = 2 > 0$ everywhere: every point is gushing. For the
swirl $(-y,x)$, $\operatorname{div}\vec F = 0$: swirling moves fluid around but never
makes it pile up or thin out.
\end{intuitionbox}

\subsection{Curl: how much a point is a whirlpool}

\begin{definition}[Curl]
\label{ch35:def-curl}
In three dimensions, $\operatorname{curl}\vec F = \grad\times\vec F$, the vector
\[
  \operatorname{curl}\vec F =
  \left(\pdv{R}{y}-\pdv{Q}{z},\;\; \pdv{P}{z}-\pdv{R}{x},\;\; \pdv{Q}{x}-\pdv{P}{y}\right).
\]
In the plane, only the last component survives; the \emph{scalar curl} of
$(P,Q)$ is $\pdv{Q}{x}-\pdv{P}{y}$ -- exactly the integrand of Green's theorem.
\end{definition}

\begin{intuitionbox}
Curl measures \emph{spin}. Drop a tiny paddlewheel into the fluid at a point. If
the flow makes it rotate, the curl is nonzero, and it points along the axle (by
the right-hand rule), with magnitude twice the angular speed of the spin. The
swirl $(-y,x)$ has scalar curl $1-(-1)=2$ everywhere -- a paddlewheel anywhere in
it spins counterclockwise. The source $(x,y)$ has curl $0-0=0$: gas rushing
radially outward never spins a paddlewheel, it just carries it straight out.
\end{intuitionbox}

\begin{trapbox}
A field can \emph{circle the origin} and still have zero curl everywhere it is
defined -- ``curls around'' (global) and ``has curl'' (local) are different. The
hole field $\big(\frac{-y}{x^2+y^2},\frac{x}{x^2+y^2}\big)$ has scalar curl $0$ at
every point of its domain, yet its loop integral around the origin is $2\pi$.
Curl is a local paddlewheel test; circulation around a big loop is a global sum.
Green's theorem links them only when the field is smooth \emph{everywhere inside}
the loop -- and this field is not (it blows up at the origin).
\end{trapbox}

\subsection{The microscopic geometry: where the formulas come from}

Why is divergence $\pdv{P}{x} + \pdv{Q}{y}$? Why is curl $\pdv{Q}{x} - \pdv{P}{y}$?
Neither formula is an arbitrary algebraic definition; both emerge directly from inspecting
a single microscopic coordinate cell $[x, x+\Delta x] \times [y, y+\Delta y]$.

\begin{figure}[H]
\centering
\begin{tikzpicture}[scale=0.88,>={Stealth[length=4pt]}]
  % Left panel: Divergence as Net Flux
  \begin{scope}[shift={(0,0)}]
    \draw[fill=red!3,draw=red!15,rounded corners=4pt] (-3.6,-2.8) rectangle (4.2,3.35);
    \node[font=\small\bfseries,anchor=north] at (0.3,3.2) {(a) Divergence: Net Outward Flux};
    \draw[fill=white,draw=gray!50,thick] (-1.4,-1.4) rectangle (1.4,1.4);
    % Right wall outflow
    \draw[->,tangentred,very thick] (1.4,0) -- ++(0.75,0) node[right=1pt,font=\scriptsize] {$P(x+\Delta x,y)\,\Delta y$};
    % Left wall inflow
    \draw[->,tangentred,very thick] (-2.15,0) -- (-1.45,0) node[pos=0.0,below=3pt,font=\scriptsize] {$P(x,y)\,\Delta y$};
    % Top wall outflow
    \draw[->,tangentred,very thick] (0,1.4) -- ++(0,0.7) node[above=1pt,font=\scriptsize] {$Q(x,y+\Delta y)\,\Delta x$};
    % Bottom wall inflow
    \draw[->,tangentred,very thick] (0,-2.15) -- (0,-1.45) node[pos=0.25,right=3pt,font=\scriptsize] {$Q(x,y)\,\Delta x$};
    % Dimensions
    \draw[gray!50] (-1.4,-1.65) -- (1.4,-1.65) node[pos=0.22,below=1pt,font=\scriptsize] {$\Delta x$};
    \draw[gray!50] (-1.65,-1.4) -- (-1.65,1.4) node[pos=0.78,left=1pt,font=\scriptsize] {$\Delta y$};
    \node[font=\scriptsize,align=center,fill=white,inner sep=2pt,rounded corners=2pt] at (0,0) {$\text{Net flux} \approx$\\$\left(\pdv{P}{x} + \pdv{Q}{y}\right)\Delta A$};
  \end{scope}

  % Right panel: Curl as Net Circulation
  \begin{scope}[shift={(8.4,0)}]
    \draw[fill=blue!3,draw=blue!15,rounded corners=4pt] (-3.6,-2.8) rectangle (4.2,3.35);
    \node[font=\small\bfseries,anchor=north] at (0.3,3.2) {(b) Curl: Net Circulation};
    \draw[fill=white,draw=gray!50,thick] (-1.4,-1.4) rectangle (1.4,1.4);
    % Bottom edge (forward in x)
    \draw[->,curveblue,very thick] (-1.2,-1.4) -- (1.2,-1.4) node[midway,below=2pt,font=\scriptsize] {$+P(x,y)\,\Delta x$};
    % Right edge (forward in y)
    \draw[->,curveblue,very thick] (1.4,-1.2) -- (1.4,1.2) node[midway,right=2pt,font=\scriptsize] {$+Q(x+\Delta x,y)\,\Delta y$};
    % Top edge (backward in x)
    \draw[->,curveblue,very thick] (1.2,1.4) -- (-1.2,1.4) node[midway,above=2pt,font=\scriptsize] {$-P(x,y+\Delta y)\,\Delta x$};
    % Left edge (backward in y)
    \draw[->,curveblue,very thick] (-1.4,1.2) -- (-1.4,-1.2) node[midway,left=2pt,font=\scriptsize] {$-Q(x,y)\,\Delta y$};
    % Center summary
    \node[font=\scriptsize,align=center,fill=white,inner sep=2pt,rounded corners=2pt] at (0,0) {$\text{Net swirl} \approx$\\$\left(\pdv{Q}{x} - \pdv{P}{y}\right)\Delta A$};
  \end{scope}
\end{tikzpicture}
\caption{The microscopic origin of divergence and curl. (a) Divergence counts the net outward
flux per unit area across the four walls of a cell: the difference in horizontal flow across $\Delta x$
gives $\pdv{P}{x}$, plus the difference in vertical flow across $\Delta y$ gives $\pdv{Q}{y}$.
(b) Curl counts the net counterclockwise circulation around the four edges: the difference between the
right and left edges gives $+\pdv{Q}{x}\,\Delta x\Delta y$, while the difference between the top and bottom
edges (running backward on top) gives $-\pdv{P}{y}\,\Delta x\Delta y$.}
\label{ch35:fig-div-curl-box}
\end{figure}

\begin{aestheticbox}
Notice the symmetry. Green's theorem in \emph{circulation form}
($\oint P\,\dd x+Q\,\dd y = \iint (Q_x-P_y)$) integrates the \emph{curl}. There is
a twin, the \emph{flux form}, that integrates the \emph{divergence}:
\[
  \oint_C \vec F\cdot \vec n\,\dd s = \iint_D \left(\pdv{P}{x}+\pdv{Q}{y}\right)\dd A
  = \iint_D \operatorname{div}\vec F\,\dd A,
\]
where $\vec n$ is the outward unit normal. The same Green's theorem, read two ways,
is the seed of \emph{both} theorems to come: extending its circulation form into
3D gives Stokes' theorem; extending its flux form into 3D gives the divergence
theorem. Both theorems generalize this single planar relationship.
\end{aestheticbox}

\section{Flux and the divergence theorem: the boundary skin}
\label{ch35:divthm}

\subsection{Flux across a surface}

In 3D the analog of ``how much field crosses the boundary loop'' becomes ``how much
field crosses the boundary \emph{surface}.'' That quantity is the \emph{flux}.

\begin{definition}[Flux integral]
\label{ch35:def-flux}
Let $S$ be a surface with a chosen unit normal direction $\vec n$. The flux of
$\vec F$ through $S$ is
\[
  \iint_S \vec F\cdot\vec n\,\dd S.
\]
Only the component of $\vec F$ \emph{piercing} the surface (along $\vec n$)
contributes; field sliding \emph{along} the surface carries nothing through.
\end{definition}

\begin{intuitionbox}
Hold a hoop in a river. The flux is the volume of water passing through the hoop
per second. Hold it face-on to the current and the flux is large; turn it edge-on
and the flux drops to zero, because now the water slides past the rim instead of
through the hoop. That ``how much is going through'' -- counting only the piercing
component $\vec F\cdot\vec n$ -- is flux.
\end{intuitionbox}

\subsection{The divergence theorem}

\begin{theorem}[Divergence theorem / Gauss's theorem]
\label{ch35:thm-div}
Let $V$ be a solid region in space bounded by a closed surface $S=\partial V$ with
outward unit normal $\vec n$. If $\vec F$ has continuous partials on $V$, then
\[
  \iiint_V \operatorname{div}\vec F\,\dd V \;=\; \oiint_{S}\vec F\cdot\vec n\,\dd S.
\]
\end{theorem}

\begin{happeningbox}
The proof is the same cancellation story, one dimension up. Chop the solid $V$ into
tiny boxes. Compute the net flux out of each tiny box -- its outward push minus its
inward pull on all six faces. Now sum over all the boxes. Every \emph{internal}
face is shared by two adjacent boxes: what flows \emph{out} of one box through that
face flows \emph{into} its neighbor, so the two flux contributions are equal and
opposite and \textbf{cancel}. Only the faces on the \emph{outer skin} of $V$
survive, with no neighbor to cancel them -- and their total is the flux through
$S=\partial V$. Meanwhile the net flux out of one tiny box equals
$\operatorname{div}\vec F\times(\text{box volume})$ -- that is literally what
divergence \emph{is}, the outflow per unit volume. Summing gives
$\iiint_V \operatorname{div}\vec F\,\dd V$. Interior fluxes cancel; the boundary
records the net source inside.
\end{happeningbox}

\begin{aestheticbox}
The divergence theorem expresses a conservation law in geometric form.
All the sources scattered throughout a solid -- every point where fluid is generated --
must account exactly for the net flux passing through its surface.
In a steady state, whatever is generated inside must pass through the enclosing surface.
This single identity underlies Gauss's law of electrostatics (charge inside equals field flux out),
the continuity equation of fluid flow, and the heat equation's energy balance.
\end{aestheticbox}

\begin{figure}[H]
\centering
\begin{tikzpicture}[scale=1.0,>={Stealth[length=4pt]}]
% a solid blob with outward normals and interior source arrows
\draw[fill=red!5,draw=red!50!black,thick] plot[smooth cycle,tension=0.8]
  coordinates {(0,1.1)(1.2,0.7)(1.4,-0.4)(0.4,-1.1)(-0.9,-0.7)(-1.2,0.5)};
% interior source arrows (divergence > 0)
\foreach \a in {0,60,...,300}{
  \draw[->,tangentred,thick] (\a:0.35) -- (\a:0.75);
}
\node[fill=white,inner sep=2pt,rounded corners=2pt,font=\scriptsize\bfseries,tangentred!80!black] at (0,0) {source};
% outward normals on boundary
\draw[->,curveblue,very thick] (1.3,0.15) -- (1.95,0.2) node[right,font=\scriptsize] {$\vec n$};
\draw[->,curveblue,very thick] (0.4,-1.1) -- (0.55,-1.75) node[below,font=\scriptsize] {$\vec n$};
\draw[->,curveblue,very thick] (-1.1,0.0) -- (-1.75,-0.05) node[left,font=\scriptsize] {$\vec n$};
\draw[->,curveblue,very thick] (0,1.1) -- (0.05,1.75) node[above,font=\scriptsize] {$\vec n$};
\node[curveblue!80!black,font=\scriptsize\bfseries] at (2.4,0.6) {outward flux};
\node[font=\small] at (0,-2.5) {$\iiint_V \operatorname{div}\vec F\,\dd V = \iint_{\partial V} \vec F\cdot\vec n\,\dd S$};
\end{tikzpicture}
\caption{The divergence theorem. Sources generating fluid in the interior (red) must
balance exactly against the net flux escaping through the enclosing skin (blue
outward normals $\vec n$). Internal flows cancel between adjoining cells; only the outer
boundary $\partial V$ records the net generation inside.}
\label{ch35:fig-div}
\end{figure}

\begin{example}[Divergence theorem on a ball]
\label{ch35:ex-div}
Let $\vec F=(x,y,z)$ and let $V$ be the ball of radius $a$, with $S$ its sphere.

\textbf{Left side.} $\operatorname{div}\vec F = 1+1+1 = 3$, so
\[
  \iiint_V 3\,\dd V = 3\cdot \text{Vol}(V) = 3\cdot\frac{4}{3}\pi a^3 = 4\pi a^3.
\]

\textbf{Right side.} On the sphere of radius $a$ the outward normal is
$\vec n = \frac{1}{a}(x,y,z)$, so
$\vec F\cdot\vec n = \frac{1}{a}(x^2+y^2+z^2) = \frac{a^2}{a} = a$, a constant.
Then
\[
  \oiint_S \vec F\cdot\vec n\,\dd S = a\cdot \text{Area}(S) = a\cdot 4\pi a^2 = 4\pi a^3.
\]
The two agree. The uniform source density $3$ filling the ball is paid off exactly
by the outward flux through the sphere.
\end{example}

\section{Stokes' theorem: the boundary loop of a surface}
\label{ch35:stokes}

Now the last of the three. Green's theorem lived in the flat plane; Stokes' theorem
lifts it into 3D, where the surface $S$ may be any curved membrane in space, like a
soap film, with the line integral run around its rim $\partial S$.

\begin{theorem}[Stokes' theorem]
\label{ch35:thm-stokes}
Let $S$ be an oriented surface in space bounded by a simple closed curve
$C=\partial S$, oriented compatibly (right-hand rule). If $\vec F$ has continuous
partials, then
\[
  \oint_C \vec F\cdot\dd\vec r \;=\; \iint_S (\operatorname{curl}\vec F)\cdot\vec n\,\dd S.
\]
\end{theorem}

\begin{intuitionbox}
The total circulation you feel marching once around the rim of a membrane equals
the sum of all the little paddlewheel spins (the curl) spread across the membrane.
And here is the wonder: \emph{any} surface with the same rim gives the same answer.
A flat disk, a tall dome, a curved sheet -- if they share a boundary loop, the
curl-flux through each is identical, because only the boundary values determine the integral. (Since
$\operatorname{div}(\operatorname{curl}\vec F)=0$ always, the divergence theorem
guarantees this independence: curl has no sources, so its flux depends only on the
boundary.) Green's theorem is Stokes' theorem for the special case where
$S$ is a flat region lying in the plane.
\end{intuitionbox}

\begin{happeningbox}
The same cancellation mechanism applies on a curved surface. Tile the membrane $S$ with tiny patches. Around each
patch run the local circulation; by Green's theorem in miniature, that circulation
equals $(\operatorname{curl}\vec F\cdot\vec n)\times(\text{patch area})$.
Sum over all patches. The shared edges between neighboring patches are traversed in
opposite directions and \textbf{cancel}; only the outer rim $C=\partial S$ survives.
Result: rim circulation $=$ total curl-flux.
\end{happeningbox}

\begin{figure}[H]
\centering
\begin{tikzpicture}[scale=0.92,>={Stealth[length=4pt]}]
  % Panel (a): Curved Cap S_1
  \begin{scope}[shift={(0,0)}]
    \draw[fill=green!3,draw=green!15,rounded corners=4pt] (-2.6,-2.3) rectangle (2.6,2.3);
    \node[font=\small\bfseries,anchor=north] at (0,2.15) {(a) Curved dome $S_1$};
    % dome surface
    \draw[fill=targetgreen!12,draw=targetgreen!60!black,thick]
      (-1.8,0) .. controls (-1.2,1.3) and (1.2,1.3) .. (1.8,0)
      .. controls (1.2,-0.4) and (-1.2,-0.4) .. (-1.8,0) -- cycle;
    % curl normal arrows poking through dome
    \foreach \x in {-1.0,-0.4,0.4,1.0}{
      \draw[->,targetgreen!80!black,thick] (\x,0.45) -- (\x,1.05);
    }
    \node[targetgreen!80!black,font=\scriptsize] at (1.1,1.15) {$\operatorname{curl}\vec F$};
    % boundary ellipse rim
    \draw[curveblue,very thick] (0,0) ellipse (1.8cm and 0.4cm);
    \draw[->,curveblue,ultra thick] (-0.2,-0.4) -- (0.3,-0.4);
    \draw[->,curveblue,ultra thick] (0.2,0.4) -- (-0.3,0.4);
    \node[curveblue,font=\scriptsize\bfseries] at (0,-0.65) {$C=\partial S_1$};
    \node[font=\scriptsize] at (0,-1.75) {$\iint_{S_1} (\operatorname{curl}\vec F)\cdot\vec n\,\dd S$};
  \end{scope}

  % Panel (b): Flat Disk S_2
  \begin{scope}[shift={(5.8,0)}]
    \draw[fill=green!3,draw=green!15,rounded corners=4pt] (-2.6,-2.3) rectangle (2.6,2.3);
    \node[font=\small\bfseries,anchor=north] at (0,2.15) {(b) Flat disk $S_2$ (same rim)};
    % flat disk
    \filldraw[fill=targetgreen!12,draw=targetgreen!60!black,thick]
      (0,0) ellipse (1.8cm and 0.4cm);
    % curl vertical arrows
    \foreach \x in {-1.0,-0.4,0.4,1.0}{
      \draw[->,targetgreen!80!black,thick] (\x,0) -- (\x,0.7);
    }
    \node[targetgreen!80!black,font=\scriptsize] at (1.1,0.85) {$\operatorname{curl}\vec F$};
    % boundary ellipse rim
    \draw[curveblue,very thick] (0,0) ellipse (1.8cm and 0.4cm);
    \draw[->,curveblue,ultra thick] (-0.2,-0.4) -- (0.3,-0.4);
    \draw[->,curveblue,ultra thick] (0.2,0.4) -- (-0.3,0.4);
    \node[curveblue,font=\scriptsize\bfseries] at (0,-0.65) {$C=\partial S_2$};
    \node[font=\scriptsize] at (0,-1.75) {$\iint_{S_2} (\operatorname{curl}\vec F)\cdot\vec n\,\dd S$};
  \end{scope}
\end{tikzpicture}
\caption{Stokes' theorem and surface independence. The circulation $\oint_C \vec F\cdot\dd\vec r$
around the boundary rim $C$ equals the total flux of $\operatorname{curl}\vec F$ through
\emph{any} surface that caps $C$. (a) A curved dome $S_1$ and (b) a flat disk $S_2$ sharing
the same rim $C=\partial S$ capture the exact same total curl flux, because any two such surfaces
together form a closed container through which the divergence-free curl field has zero net flux.}
\label{ch35:fig-stokes}
\end{figure}

\begin{example}[Stokes on a hemisphere]
\label{ch35:ex-stokes}
Let $\vec F=(-y,x,0)$ (the swirl, lifted into 3D) and let $S$ be the upper unit
hemisphere $x^2+y^2+z^2=1$, $z\ge0$, oriented upward. Its boundary $C$ is the unit
circle in the $xy$-plane.

\textbf{Surface side.}
$\operatorname{curl}\vec F = (0,0,\,Q_x-P_y)=(0,0,1-(-1))=(0,0,2)$.
The simplest membrane with the same rim is the flat disk $z=0$, where
$\vec n=(0,0,1)$, so $(\operatorname{curl}\vec F)\cdot\vec n = 2$ and
\[
  \iint_S (\operatorname{curl}\vec F)\cdot\vec n\,\dd S = 2\cdot(\text{area of unit disk}) = 2\pi.
\]
(We are allowed to swap the hemisphere for the flat disk precisely because Stokes
says only the shared rim matters.)

\textbf{Rim side.} Parametrize $C$ by $x=\cos t$, $y=\sin t$, $z=0$:
\[
  \oint_C \vec F\cdot\dd\vec r = \int_0^{2\pi}\big[(-\sin t)(-\sin t)+(\cos t)(\cos t)\big]\dd t = 2\pi.
\]
Both give $2\pi$. The curl pointing up through the cap is precisely the swirl you
feel walking the equator.
\end{example}

\begin{historybox}
The story of vector calculus begins not in a university lecture hall, but in a grain
windmill in Nottingham. In 1828, George Green, a self-taught miller who had only one
year of formal schooling as a child, published a private essay on electricity and
magnetism. Green was trying to solve a concrete puzzle: if electrical charges or
heat sources are buried invisibly inside a solid body, how can you deduce what they
are doing using only measurements taken by meters placed on the outside skin?

At the time, mathematicians worked in coordinate components. They wrote
out separate formulas for the $x$, $y$, and $z$ directions. Modern vector
notation did not exist: neither dot products nor cross products, nor the operational symbols
$\grad$, $\operatorname{div}$, and $\operatorname{curl}$, had been invented yet. Green
realized that the sum of microscopic changes inside an enclosure must cancel along
every interior boundary, leaving only a reckoning at the outer edge.

Because Green was an unknown outsider, his self-funded booklet was sold to only 51
subscribers, mostly local friends, and was completely ignored by the scientific
world. It was almost lost forever. In 1845, a young Scottish physicist named William
Thomson (later Lord Kelvin) stumbled upon a surviving copy, immediately saw its
importance, and circulated it across Europe. Five years later, in 1850, Thomson wrote a
letter to his friend George Gabriel Stokes containing the three-dimensional version
of the boundary principle on a curved surface. Stokes never published a proof; instead,
he set it as Question~8 on the 1854 Cambridge Smith's Prize examination, an advanced
mathematics competition. Sitting that very exam was an undergraduate named James Clerk
Maxwell.

Maxwell seized the theorem and made curl and divergence central to his
equations of electromagnetism, showing for the first time that light is an
electromagnetic wave. A lively debate soon followed over how to formulate these
ideas. Peter Guthrie Tait championed four-dimensional numbers called quaternions,
debating with Josiah Willard Gibbs and Oliver Heaviside,
who argued for the three-dimensional vectors used today.
By 1900, Henri Poincar\'e and \'Elie Cartan unified these results into
the generalized Stokes theorem $\int_{\partial M}\omega = \int_M \dd\omega$,
where $\omega$ is a differential form (a tidy geometric package ready to be integrated)
and $\partial M$ is the boundary of a smooth region $M$.

Yet even with this beautiful boundary geometry in hand, the fundamental physics of
fluid flow remains an open frontier. The Navier--Stokes existence and smoothness
problem -- one of the million-dollar Millennium Prize Problems -- asks whether the
three-dimensional vector equations governing flowing water and air always have smooth,
predictable solutions for all time, or whether the curl of the velocity field (its
local whirlpool spin) can spontaneously concentrate and blow up to infinity in finite
time.
\end{historybox}

\section{One theorem, four faces}
\label{ch35:unification}

Step back and see the family resemblance. In each statement, an integral of a
\emph{derivative} over a region equals an integral of the original \emph{field}
over the region's \emph{boundary}:

\begin{center}
\renewcommand{\arraystretch}{1.6}
\begin{tabular}{@{}lll@{}}
\toprule
\textbf{Theorem} & \textbf{Inside (a derivative)} & \textbf{Boundary (the field)} \\
\midrule
FTC & $\displaystyle\int_a^b F'(x)\,\dd x$ & $F(b)-F(a)$ \\
Gradient thm & $\displaystyle\int_C \grad f\cdot\dd\vec r$ & $f(B)-f(A)$ \\
Green & $\displaystyle\iint_D (Q_x-P_y)\,\dd A$ & $\displaystyle\oint_{\partial D}\vec F\cdot\dd\vec r$ \\
Stokes & $\displaystyle\iint_S \operatorname{curl}\vec F\cdot\vec n\,\dd S$ & $\displaystyle\oint_{\partial S}\vec F\cdot\dd\vec r$ \\
Divergence & $\displaystyle\iiint_V \operatorname{div}\vec F\,\dd V$ & $\displaystyle\oiint_{\partial V}\vec F\cdot\vec n\,\dd S$ \\
\bottomrule
\end{tabular}
\end{center}

\begin{aestheticbox}
The boundary operator $\partial$ lowers dimension by one: the boundary of a solid is
a surface, of a surface a curve, of a curve two endpoints. The derivative operator
($\grad$, curl, div) is the perfect partner: applied inside, it produces exactly
what, when summed, the boundary can account for. The two are \emph{adjoint} -- that
is the one-sentence content of the generalized Stokes theorem,
$\int_M \dd\omega = \int_{\partial M}\omega$. Every theorem you have met in this
book, from $\int F' = F(b)-F(a)$ onward, is a special case of that single formula.
\end{aestheticbox}

\begin{connectionbox}
This unification provides a compact mathematical language for physical field theories. Maxwell's
equations say $\operatorname{div}\vec E = \rho/\eps_0$ (charges are sources of the
electric field) and $\operatorname{curl}\vec B = \mu_0\vec J + \cdots$ (currents
swirl the magnetic field). Feed these into the divergence and Stokes theorems and
out come Gauss's law and Ampere's law in their boundary form -- the form you can
actually measure with a wire and a meter. The local differential law and the global
integral law are the same statement seen from inside and from the boundary.
\end{connectionbox}

\begin{labbox}
Verify the unification on a concrete example. Take $\vec F=(-y,x)$.
(1) Compute its scalar curl. (2) Compute its divergence. (3) Use Green's theorem to
find its circulation around the boundary of the square $[0,1]\times[0,1]$, then
check it by integrating around the four edges directly. (4) Now compute its
\emph{flux} out of that same square using the flux form of Green's theorem -- and
explain, before you compute, why you expect $0$. (The swirl spins but does not
spread.)
\end{labbox}

\section{The Bridge to Complex Analysis: Analytic Functions, Cauchy--Riemann, and Contour Integrals}
\label{ch35:bridge-complex-analysis}

Throughout this chapter, Green's theorem, Stokes' theorem, and the divergence theorem established a single unifying principle: an integral of a derivative over an interior region is completely recorded by the field along the boundary. While historically developed to describe fluid vortices and electromagnetic fields, Green's theorem produces an even more striking payoff in two dimensions when the coordinate plane $\R^2$ is identified with the complex plane $\C$ via $z = x + iy$.

Here, Green's theorem becomes the foundation of \emph{Complex Analysis}---a discipline where functions possessing a single complex derivative automatically possess infinitely many, and where difficult definite integrals are evaluated by counting how paths wind around isolated singularities. A \emph{singularity} is simply a point where a formula fails to be well-behaved, such as a location where it attempts to divide by zero.

\subsection{Complex Differentiability and the Amplitwist}

A complex function $f: \C \to \C$ maps an input point $z = x + iy$ to an output $w = u(x,y) + iv(x,y)$, where $u, v: \R^2 \to \R$ are its real and imaginary parts. 

Just as in single-variable calculus, the complex derivative of $f$ at a point $z_0$ is defined as the limit of a difference quotient:
\begin{equation}
f'(z_0) = \lim_{\Delta z \to 0} \frac{f(z_0 + \Delta z) - f(z_0)}{\Delta z}.
\label{ch35:eq-complex-derivative}
\end{equation}
In single-variable real calculus, the increment $\Delta x$ can approach zero from only two directions: left or right. In the complex plane, however, $\Delta z = \Delta x + i\Delta y$ can approach zero along an infinite family of two-dimensional trajectories. For the complex derivative $f'(z_0)$ to exist, the quotient must converge to the exact same complex number regardless of the trajectory of approach.

Compare two perpendicular directions of approach:
\begin{enumerate}
\item \textbf{Horizontal approach} ($\Delta z = \Delta x \in \R$, so $\Delta y = 0$):
\[
f'(z) = \lim_{\Delta x \to 0} \frac{u(x+\Delta x, y) - u(x,y)}{\Delta x} + i \lim_{\Delta x \to 0} \frac{v(x+\Delta x, y) - v(x,y)}{\Delta x} = u_x + i v_x.
\]
\item \textbf{Vertical approach} ($\Delta z = i\Delta y$ with $\Delta y \in \R$, so $\Delta x = 0$):
\[
f'(z) = \lim_{\Delta y \to 0} \frac{u(x, y+\Delta y) - u(x,y)}{i\Delta y} + i \lim_{\Delta y \to 0} \frac{v(x, y+\Delta y) - v(x,y)}{i\Delta y} = \frac{u_y + i v_y}{i}.
\]
Multiplying numerator and denominator by $-i$ (recalling that $1/i = -i$):
\[
\frac{u_y + i v_y}{i} = -i(u_y + i v_y) = v_y - i u_y.
\]
\end{enumerate}

For $f'(z)$ to exist, these two directional limits must be identical:
\[
u_x + i v_x = v_y - i u_y.
\]
Equating real and imaginary components yields the celebrated \textbf{Cauchy--Riemann equations}:
\begin{equation}
u_x = v_y \qquad \text{and} \qquad u_y = -v_x.
\label{ch35:eq-cauchy-riemann}
\end{equation}

\subsection{The Geometric View: Pure Rotation and Scaling}

What do the Cauchy--Riemann equations mean geometrically?

View $f$ as a multivariable mapping from the plane $\R^2$ to $\R^2$, sending $(x,y)$ to $(u(x,y), v(x,y))$. Its local linear approximation is governed by the $2 \times 2$ Jacobian matrix:
\[
J = \begin{pmatrix} u_x & u_y \\ v_x & v_y \end{pmatrix}.
\]
Imposing the Cauchy--Riemann equations $v_y = u_x$ and $u_y = -v_x$, the Jacobian specializes to:
\[
J = \begin{pmatrix} u_x & -v_x \\ v_x & u_x \end{pmatrix}.
\]
Let $r = \sqrt{u_x^2 + v_x^2} = \abs{f'(z)}$, and let $\theta = \operatorname{Arg}(u_x + i v_x) = \operatorname{Arg}(f'(z))$. Then this matrix factors into:
\[
J = r \begin{pmatrix} \cos\theta & -\sin\theta \\ \sin\theta & \cos\theta \end{pmatrix}.
\]
The Jacobian is not an arbitrary linear transformation: it is a pure counterclockwise rotation by angle $\theta$ accompanied by an isotropic scaling by factor $r$.

As Tristan Needham termed it in \emph{Visual Complex Analysis}, a complex derivative acts locally as an \textbf{amplitwist}: it \emph{amplifies} lengths by $r = \abs{f'(z)}$ and \emph{twists} directions by angle $\theta = \operatorname{Arg}(f'(z))$. Because all directions are scaled equally and angles between vectors are preserved, every complex-differentiable (or \emph{holomorphic}) map is \emph{conformal}: tiny squares in the domain map to tiny, rotated squares in the range, free of any shear or directional stretching.

\begin{figure}[H]
\centering
\begin{tikzpicture}[scale=0.92,>={Stealth[length=5pt]}]
  % Panel A: Amplitwist
  \begin{scope}[shift={(0,0)}]
    \draw[fill=blue!3,draw=blue!15,rounded corners=4pt] (-3.8,-3.2) rectangle (3.8,3.2);
    \node[font=\small\bfseries,anchor=north] at (0,3.05) {(a) Amplitwist: local conformal grid mapping};
    
    % Domain z-plane (left)
    \node[font=\scriptsize\bfseries, anchor=north] at (-2.1, 2.5) {$z$-plane};
    \draw[->, gray!60] (-3.3,-0.8) -- (-0.9,-0.8) node[right, font=\tiny] {$x$};
    \draw[->, gray!60] (-2.1,-1.8) -- (-2.1,1.8) node[above, font=\tiny] {$y$};
    % Cartesian grid of small squares
    \foreach \gx in {-2.9,-2.5,-2.1,-1.7,-1.3}{
      \draw[gray!30, thin] (\gx, -0.6) -- (\gx, 1.0);
    }
    \foreach \gy in {-0.6,-0.2,0.2,0.6,1.0}{
      \draw[gray!30, thin] (-2.9, \gy) -- (-1.3, \gy);
    }
    % Highlighted square in domain
    \filldraw[fill=curveblue!20, draw=curveblue, thick] (-2.1, 0.2) rectangle (-1.7, 0.6);
    \node[font=\tiny, curveblue!80!black] at (-1.9, 0.4) {$\dd z$};
    
    % Mapping arrow
    \draw[->, quotientorange, line width=1.8pt] (-0.7, 0.2) -- (0.5, 0.2) node[midway, above=1pt, font=\scriptsize\bfseries] {$f(z)$};
    \node[font=\tiny, align=center] at (-0.1, -0.25) {amplitwist\\scale $r$, rot $\theta$};
    
    % Target w-plane (right)
    \node[font=\scriptsize\bfseries, anchor=north] at (2.1, 2.5) {$w$-plane};
    \draw[->, gray!60] (0.9,-0.8) -- (3.3,-0.8) node[right, font=\tiny] {$u$};
    \draw[->, gray!60] (2.1,-1.8) -- (2.1,1.8) node[above, font=\tiny] {$v$};
    
    % Rotated grid in target
    \begin{scope}[shift={(2.35, 0.2)}, rotate=32, scale=1.15]
      \foreach \gx in {-0.8,-0.4,0.0,0.4,0.8}{
        \draw[targetgreen!30, thin] (\gx, -0.8) -- (\gx, 0.8);
      }
      \foreach \gy in {-0.8,-0.4,0.0,0.4,0.8}{
        \draw[targetgreen!30, thin] (-0.8, \gy) -- (0.8, \gy);
      }
      % Highlighted rotated square
      \fill[white] (0.0, 0.0) rectangle (0.4, 0.4);
      \filldraw[fill=targetgreen!25, draw=targetgreen, thick] (0.0, 0.0) rectangle (0.4, 0.4);
      \node[font=\tiny, targetgreen!80!black] at (0.2, 0.2) {$\dd w$};
    \end{scope}
    
    \node[font=\tiny, align=center] at (0,-2.6) {Jacobian $J = r R_\theta$: rotation $+$ isotropic scaling.\\Right angles ($90^\circ$) and tiny squares preserved.};
  \end{scope}

  % Panel B: Contour Deformation Around a Pole
  \begin{scope}[shift={(8.4,0)}]
    \draw[fill=blue!3,draw=blue!15,rounded corners=4pt] (-3.8,-3.2) rectangle (3.8,3.2);
    \node[font=\small\bfseries,anchor=north] at (0,3.05) {(b) Contour deformation around pole $z_0$};
    
    % Outer loop C
    \draw[curveblue, line width=1.5pt] (-2.6,-1.2) to[out=70,in=190] (-0.5,2.0) to[out=10,in=110] (2.5,1.2) to[out=-70,in=10] (1.8,-1.8) to[out=190,in=-110] (-2.6,-1.2);
    \draw[->, curveblue, ultra thick] (-0.8,1.98) -- (-0.4,2.02) node[above=2pt, font=\scriptsize] {$C$};
    \draw[->, curveblue, ultra thick] (2.2,0.2) -- (2.0,-0.4);
    
    % Pole z_0
    \coordinate (Z0) at (0.0, 0.0);
    \filldraw[tangentred] (Z0) circle (2.5pt) node[below left=2pt, font=\scriptsize\bfseries, text=black] {$z_0$ (pole)};
    
    % Small circular contour C_eps around z_0
    \draw[targetgreen, line width=1.5pt] (Z0) circle (0.85cm);
    \draw[->, targetgreen, ultra thick] (0.85, 0.0) arc (0:80:0.85cm) node[above right=-2pt, font=\scriptsize] {$C_\eps$};
    
    % Connecting cut corridor with opposing cancellation arrows
    \draw[gray!50, thick, dashed] (0.85, 0.08) -- (2.35, 0.08);
    \draw[gray!50, thick, dashed] (0.85, -0.08) -- (2.35, -0.08);
    \draw[->, gray!80!black, thick] (1.3, 0.18) -- (1.8, 0.18);
    \draw[->, gray!80!black, thick] (1.8, -0.18) -- (1.3, -0.18);
    \node[font=\tiny, gray!70!black, above] at (1.6, 0.25) {opposing cut cancels};
    
    % Region D label
    \node[font=\tiny, gray!80!black] at (-1.5, 0.9) {$D$ ($\operatorname{curl}=0$)};
    
    \node[font=\tiny, align=center] at (0,-2.6) {Green's thm on $D$: $\oint_C \frac{\dd z}{z-z_0} = \oint_{C_\eps} \frac{\dd z}{z-z_0} = 2\pi i$.\\Winding number counts enclosed poles.};
  \end{scope}
\end{tikzpicture}
\caption{The bridge from Green's theorem to complex analysis. (a) A complex function $f(z)$ satisfies the Cauchy--Riemann equations if and only if its derivative acts locally as an amplitwist: rotating by $\theta = \operatorname{Arg}(f')$ and scaling isotropically by $r = \abs{f'}$, preserving infinitesimal squares. (b) Green's theorem implies Cauchy's integral theorem ($\oint_C f(z)\,\dd z = 0$ on holomorphic domains). When an isolated pole $z_0$ is enclosed, Green's theorem on the annular region between $C$ and a small circle $C_\eps$ proves the contour can be deformed onto $C_\eps$, evaluating to $2\pi i$.}
\label{ch35:fig-cauchy-contour}
\end{figure}

\subsection{Cauchy's Integral Theorem in Three Lines via Green's Theorem}

Now consider the contour integral of a complex function $f(z)$ around a closed curve $C = \partial D$ bounding a simply connected region $D \subset \C$. 

Expand the complex integrand and differential displacement into real and imaginary components:
\[
f(z) = u + iv, \qquad \dd z = \dd x + i\,\dd y.
\]
Multiplying these expressions:
\[
f(z)\,\dd z = (u + iv)(\dd x + i\,\dd y) = (u\,\dd x - v\,\dd y) + i\,(v\,\dd x + u\,\dd y).
\]
Integrating around the closed boundary loop $C$ splits the complex contour integral into two ordinary real line integrals:
\begin{equation}
\oint_C f(z)\,\dd z = \oint_C (u\,\dd x - v\,\dd y) + i \oint_C (v\,\dd x + u\,\dd y).
\label{ch35:eq-contour-split}
\end{equation}

Now apply Green's Theorem, $\oint_C (P\,\dd x + Q\,\dd y) = \iint_D (Q_x - P_y)\,\dd A$, directly to both integrals:
\begin{enumerate}
\item \textbf{Real part} ($P = u$, $Q = -v$):
\[
\pdv{Q}{x} - \pdv{P}{y} = \pdv{(-v)}{x} - \pdv{u}{y} = -v_x - u_y.
\]
By the second Cauchy--Riemann equation, $u_y = -v_x$. Substituting this gives:
\[
-v_x - (-v_x) = 0.
\]
\item \textbf{Imaginary part} ($P = v$, $Q = u$):
\[
\pdv{Q}{x} - \pdv{P}{y} = \pdv{u}{x} - \pdv{v}{y} = u_x - v_y.
\]
By the first Cauchy--Riemann equation, $u_x = v_y$. Substituting this gives:
\[
v_y - v_y = 0.
\]
\end{enumerate}

Both interior curl integrands vanish identically throughout the domain $D$:
\begin{equation}
\oint_C f(z)\,\dd z = \iint_D 0\,\dd A + i \iint_D 0\,\dd A = 0.
\label{ch35:eq-cauchy-theorem}
\end{equation}

This is \textbf{Cauchy's Integral Theorem}:
\begin{theorem}[Cauchy's Integral Theorem]
If $f(z)$ is complex-differentiable (holomorphic) on and inside a simple closed contour $C$, then:
\[
\oint_C f(z)\,\dd z = 0.
\]
\end{theorem}
Cauchy's integral theorem is Green's theorem applied to complex functions: the Cauchy--Riemann equations ensure that the vector fields $(u, -v)$ and $(v, u)$ are identically curl-free.

\subsection{The Single Punctured Exception: \texorpdfstring{$\oint_{\abs{z}=R} \frac{1}{z}\,\dd z = 2\pi i$}{The Contour Integral of 1/z}}

What occurs when $f(z)$ fails to be differentiable at an isolated singularity inside $C$?

The standard test case is $f(z) = 1/z$, which is holomorphic everywhere except at the origin $z = 0$, where it is undefined. An isolated singularity where $|f(z)| \to \infty$ is called a \emph{pole}. 

Let $C$ be the circle $\abs{z} = R$ traversed counterclockwise. Parametrize $C$ in polar form:
\[
z(\theta) = R e^{i\theta}, \qquad \dd z = i R e^{i\theta}\,\dd\theta, \qquad \theta \in [0, 2\pi].
\]
Substitute this parametrization into the contour integral:
\begin{equation}
\oint_{\abs{z}=R} \frac{1}{z}\,\dd z = \int_0^{2\pi} \frac{1}{R e^{i\theta}} \left(i R e^{i\theta}\,\dd\theta\right) = \int_0^{2\pi} i\,\dd\theta = 2\pi i.
\label{ch35:eq-1-over-z}
\end{equation}

Observe the cancellation: the radius $R$ cancels completely from the calculation. Whether the circular path has radius $R = 10^5$ or $R = 10^{-5}$, the line integral is invariantly $2\pi i$. 

By applying Green's theorem to the annular region between an arbitrary smooth loop $C$ encircling the origin and a small circle $C_\eps$ of radius $\eps$, the curl is zero in the intermediate space. The connecting corridor has opposing paths that cancel (Figure~\ref{ch35:fig-cauchy-contour}b), proving that any closed contour enclosing the origin can be deformed onto $C_\eps$ without altering the integral:
\[
\oint_C \frac{1}{z}\,\dd z = \oint_{C_\eps} \frac{1}{z}\,\dd z = 2\pi i.
\]
The integral $\oint_C \frac{1}{z}\,\dd z$ acts as an integer winding counter: it evaluates to $2\pi i$ times the number of times the path $C$ winds around the singularity.

In contrast, for every other integer power $z^n$ with $n \neq -1$:
\[
\oint_{\abs{z}=R} z^n\,\dd z = \int_0^{2\pi} R^n e^{i n \theta} \left(i R e^{i\theta}\,\dd\theta\right) = i R^{n+1} \int_0^{2\pi} e^{i(n+1)\theta}\,\dd\theta = 0,
\]
because integrating a periodic sinusoidal function over full periods yields zero. The term $z^{-1}$ is the unique survivor.

\subsection{Cauchy's Residue Theorem and Signal Processing}

Any function with isolated singularities can be expanded around a pole $z_0$ in a Laurent series:
\[
f(z) = \sum_{n=-\infty}^\infty a_n (z - z_0)^n = \dots + \frac{a_{-2}}{(z-z_0)^2} + \frac{a_{-1}}{z-z_0} + a_0 + a_1(z-z_0) + \dots
\]
A \emph{Laurent series} is simply a power series that permits negative integer powers,
capturing the blow-up near the pole. When we integrate $f(z)$ around a closed loop
enclosing $z_0$, Green's theorem ensures that every positive power and every higher
negative power integrates to zero. Only the term $a_{-1}(z - z_0)^{-1}$ survives,
yielding $2\pi i a_{-1}$. The coefficient $a_{-1}$ is called the \textbf{residue} of
$f$ at $z_0$, written $\operatorname{Res}(f, z_0)$ -- it is the only nonvanishing
term in the loop integral.

Deforming a general contour around multiple enclosed isolated poles $z_1, z_2, \dots, z_k$ yields \textbf{Cauchy's Residue Theorem}:
\begin{equation}
\oint_C f(z)\,\dd z = 2\pi i \sum_{k} \operatorname{Res}(f, z_k).
\label{ch35:eq-residue-theorem}
\end{equation}

In modern engineering and physics, this residue calculation is the primary mechanism for solving linear differential equations:
\begin{itemize}
\item \textbf{Continuous-time systems and the Laplace transform}: When modeling an electrical circuit or a vibrating spring, physical laws yield differential equations relating a function to its time derivatives. The Laplace transform turns these differential equations into ordinary algebraic equations in terms of a complex frequency variable $s$. Solving for the system yields an algebraic transfer function $F(s)$. To convert the solution back into real-world time $f(t)$, one computes a contour integral in the complex plane (the Bromwich integral). Cauchy's Residue Theorem evaluates this integral directly: the sum of the residues at the poles of $F(s)$ reveals the system's natural frequencies and exponential decay rates.
\item \textbf{Discrete-time systems and digital filters}: In digital signal processing (DSP), where signals are sampled sequences of numbers $x[n]$ rather than continuous curves, the $Z$-transform plays the same role. A closed contour integral in the complex plane extracts the digital filter's time response directly from the residues of its transfer function $X(z)$.
\end{itemize}
Green's theorem, originally designed to tally fluid circulation and boundary flux, provides the foundation for evaluating these contour integrals.

\section{Looking forward: from static vector fields to differential equations}
\label{ch35:bridge-ode}

Throughout this chapter, we have treated vector fields from a \emph{global, spatial}
standpoint. We stood outside the plane or solid, viewing $\vec F(x,y)$ as a fixed
pattern of arrows painted across space, and asked accounting questions over regions:
\begin{itemize}
\item How much does the field circulate around a boundary fence $C$?
  ($\oint_C \vec F \cdot \dd\vec r$)
\item How much fluid leaks out across the enclosing skin of a solid?
  ($\oiint_S \vec F \cdot \vec n\,\dd S$)
\end{itemize}
Green's, Stokes', and the divergence theorems revealed that all microscopic curls
and sources inside cancel in pairs, leaving only a reckoning at the outer perimeter.
Even in our excursion into complex analysis, Cauchy's theorem and the Residue
Theorem operated by evaluating closed loops around singularities to decode the
vibrations of physical systems.

Now consider the trajectory of an individual point. Instead of computing
flux across an entire surface, place a particle in the flow at
an initial point $\vec r_0 = (x_0, y_0)$ and track its motion over time.

At every location $(x,y)$, the vector field plants a velocity arrow:
\[
  \vec F(x,y) = \big(P(x,y),\, Q(x,y)\big).
\]
If the vector field represents a fluid current, the speck must move in the direction
the local arrow points, with speed equal to the arrow's length. Its instantaneous
velocity $\dv{\vec r}{t}$ is dictated at every moment by its current position:
\begin{equation}
  \dv{\vec r}{t} = \vec F\big(\vec r(t)\big)
  \quad\iff\quad
  \begin{cases}
    \displaystyle \dv{x}{t} = P(x, y),\\[6pt]
    \displaystyle \dv{y}{t} = Q(x, y).
  \end{cases}
  \label{ch35:eq-field-ode}
\end{equation}
This is our first encounter with an \textbf{ordinary differential equation} (ODE):
a rule of change relating a quantity's rate of change to its current state.

\begin{figure}[H]
\centering
\begin{tikzpicture}[scale=0.92,>={Stealth[length=4pt]}]
  % Panel (a): Vector Calculus View
  \begin{scope}[shift={(0,0)}]
    \draw[fill=blue!3,draw=blue!15,rounded corners=4pt] (-3.4,-2.8) rectangle (3.4,2.8);
    \node[font=\small\bfseries,anchor=north] at (0,2.65) {(a) Vector Calculus: Global Accounting};
    % Background field arrows (swirl)
    \foreach \x in {-2.2,-1.1,0.0,1.1,2.2}{
      \foreach \y in {-1.8,-0.9,0.0,0.9,1.8}{
        \ifdim\x pt=0pt \ifdim\y pt=0pt \else
          \draw[->,gray!35,thin] (\x,\y) -- ++({-\y*0.22},{(\x+0.1)*0.22});
        \fi\else
          \draw[->,gray!35,thin] (\x,\y) -- ++({-\y*0.22},{(\x+0.1)*0.22});
        \fi
      }
    }
    % Region D with boundary C
    \filldraw[fill=curveblue!15,draw=curveblue,very thick]
      (0,0) ellipse (1.9cm and 1.2cm);
    \node[curveblue!80!black,font=\small\bfseries] at (0,0) {$D$};
    % Boundary direction arrows
    \draw[->,curveblue,ultra thick] (0,-1.2) -- (0.4,-1.2);
    \draw[->,curveblue,ultra thick] (1.9,0) -- (1.9,0.3);
    \draw[->,curveblue,ultra thick] (0,1.2) -- (-0.4,1.2);
    \draw[->,curveblue,ultra thick] (-1.9,0) -- (-1.9,-0.3);
    \node[curveblue,font=\scriptsize\bfseries] at (1.5,1.25) {$C = \partial D$};
    \node[font=\scriptsize,align=center] at (0,-2.35) {$\oint_{\partial D} \vec F\cdot\dd\vec r = \iint_D (\operatorname{curl}\vec F)\,\dd A$};
  \end{scope}

  % Panel (b): Differential Equations View
  \begin{scope}[shift={(7.3,0)}]
    \draw[fill=blue!3,draw=blue!15,rounded corners=4pt] (-3.4,-2.8) rectangle (3.4,2.8);
    \node[font=\small\bfseries,anchor=north] at (0,2.65) {(b) Differential Equations: Following Arrows};
    % Background field arrows (same swirl)
    \foreach \x in {-2.2,-1.1,0.0,1.1,2.2}{
      \foreach \y in {-1.8,-0.9,0.0,0.9,1.8}{
        \draw[->,gray!35,thin] (\x,\y) -- ++({-\y*0.22},{(\x+0.1)*0.22});
      }
    }
    % Trajectory streamline
    \draw[tangentred,very thick,->] (1.8,0) arc (0:190:1.8cm and 1.25cm);
    \draw[tangentred,very thick] (-1.77,-0.2) arc (190:350:1.8cm and 1.25cm);
    % Initial condition point
    \filldraw[black] (1.8,0) circle (2pt) node[right=2pt,font=\scriptsize\bfseries] {$\vec r(0)=\vec r_0$};
    % Velocity tangent vector at a point along the curve
    \coordinate (P) at (0,1.25);
    \filldraw[tangentred] (P) circle (1.8pt);
    \draw[->,tangentred!90!black,ultra thick] (P) -- ++(-1.1,0)
      node[above=2pt,fill=blue!3,inner sep=1pt,font=\scriptsize\bfseries] {$\dv{\vec r}{t} = \vec F(\vec r(t))$};
    \node[font=\scriptsize,align=center] at (0,-2.35) {Local steering law $\implies$ trajectory $\vec r(t)$};
  \end{scope}
\end{tikzpicture}
\caption{From vector calculus to differential equations.
(a) Vector calculus takes a global, spatial view: given a static vector field $\vec F$,
it evaluates macroscopic circulation $\oint_{\partial D} \vec F\cdot\dd\vec r$ or outward flux
by summing local derivatives across the interior region $D$.
(b) Differential equations takes a dynamical, temporal view: place a particle at an initial
state $\vec r_0$, and the local field arrow dictates its instantaneous velocity
$\dv{\vec r}{t} = \vec F(\vec r(t))$ at every moment. The resulting solution curve $\vec r(t)$
(streamline) traces the future of the system. Chapter~36 begins the study of these local laws
of change.}
\label{ch35:fig-field-to-ode}
\end{figure}

Notice the shift in perspective between the two panels of Figure~\ref{ch35:fig-field-to-ode}:
\begin{itemize}
\item In \textbf{vector calculus}, the vector field is a static canvas. We ask what
happens when we integrate across it or around closed boundary loops.
\item In \textbf{differential equations}, the vector field is a \emph{dynamical rule of motion}.
It acts as a compass and speedometer at every point in the plane. An initial point
$\vec r(0) = \vec r_0$ launches a traveler along a path $\vec r(t)$ called a
\emph{streamline} or \emph{solution trajectory}.
\end{itemize}

The geometric properties of the vector field determine the behavior of these trajectories.
If the field is conservative ($\vec F = -\grad U$), every trajectory moves downhill
along lines of steepest descent toward local minima of the potential energy $U$.
If the field is divergence-free ($\operatorname{div}\vec F = 0$), any cloud of initial
points flowing along the trajectories preserves its area for all $t$
(a fundamental principle known as Liouville's theorem in statistical mechanics).

In Chapter~36, we begin the study of differential equations from the ground up.
We will step back from multi-dimensional vectors for a moment to look at the
cleanest possible case: a single quantity $y(t)$ whose rate of change $\dv{y}{t}$
is linked directly to its own value. We will explore what a differential equation
actually is, how to visualize direction fields without solving anything, and why
these local rules of change govern phenomena ranging from radioactive decay and cooling
coffee to the orbits of planets.

\section{Exercises}
\label{ch35:exercises}

\begin{exercise}[Mental check: curl and loops]
\label{ch35:exr-mental-curl}
Suppose a vector field $\vec F$ satisfies $\operatorname{curl}\vec F = \vec 0$ everywhere in $\R^3$. Without doing any calculation, what is the line integral $\oint_C \vec F\cdot\dd\vec r$ around any simple closed loop $C$? State your reasoning in one sentence using Stokes' theorem.
\end{exercise}

\begin{exercise}[Mental check: divergence and closed containers]
\label{ch35:exr-mental-flux}
If a solid watermelon contains no fluid sources or sinks ($\operatorname{div}\vec F = 0$ everywhere inside), can more fluid flow out through its rind than flows in? What does the divergence theorem say about the net outward flux $\oiint_S \vec F\cdot\vec n\,\dd S$?
\end{exercise}

\begin{exercise}[Mental check: constant wind]
\label{ch35:exr-mental-constant}
A steady wind blows horizontally across the plains with constant velocity $\vec F(x,y) = (4, 7)$ meters per second. What is its scalar curl? What is its divergence? If you walk once around a circular running track in this wind, what is the net work done on you by the air?
\end{exercise}

\begin{exercise}
\label{ch35:exr-lineint}
Compute $\displaystyle\int_C \vec F\cdot\dd\vec r$ for $\vec F=(x+y,\,y)$ along the
straight segment from $(0,0)$ to $(3,1)$. Parametrize by $\vec r(t)=(3t,t)$,
$t\in[0,1]$, and reduce to a single integral in $t$.
\end{exercise}

\begin{exercise}
\label{ch35:exr-conservative}
Show that $\vec F=(3x^2y + e^y,\; x^3 + xe^y + 2)$ is conservative by checking
$P_y=Q_x$, then find a potential $f$. Use it to evaluate $\int_C\vec F\cdot\dd\vec r$
from $(0,0)$ to $(1,1)$ along \emph{any} path.
\end{exercise}

\begin{exercise}
\label{ch35:exr-notconservative}
For $\vec F=(y,\,-x)$, compute $\oint_C\vec F\cdot\dd\vec r$ around the unit circle
counterclockwise two ways: directly by parametrization, and via Green's theorem.
Confirm they agree, and explain why the nonzero answer proves $\vec F$ is not
conservative.
\end{exercise}

\begin{exercise}
\label{ch35:exr-green-area}
Use the area formula $\text{Area}=\tfrac12\oint_C x\,\dd y - y\,\dd x$ to find the
area enclosed by the astroid $x=\cos^3 t$, $y=\sin^3 t$, $t\in[0,2\pi]$. (You should
get $\tfrac{3}{8}\pi$.)
\end{exercise}

\begin{exercise}
\label{ch35:exr-divcurl}
For $\vec F=(x^2 y,\; -xy^2,\; z)$ compute $\operatorname{div}\vec F$ and
$\operatorname{curl}\vec F$ as functions of $(x,y,z)$. Identify whether there is any point where the
field is locally incompressible ($\operatorname{div}=0$), or show that $\operatorname{div}\vec F = 1$ everywhere.
\end{exercise}

\begin{exercise}
\label{ch35:exr-greenverify}
Verify Green's theorem for $\vec F=(xy,\,x^2)$ over the triangle with vertices
$(0,0)$, $(1,0)$, $(1,1)$, by computing both the boundary line integral (three
edges) and the double integral of $Q_x-P_y$ over the triangle.
\end{exercise}

\begin{exercise}
\label{ch35:exr-divthm}
Use the divergence theorem to compute the flux of $\vec F=(x^3,\,y^3,\,z^3)$ out of
the unit sphere. (Hint: $\operatorname{div}\vec F = 3(x^2+y^2+z^2)$; convert the
resulting triple integral to spherical coordinates.)
\end{exercise}

\begin{exercise}
\label{ch35:exr-stokes}
Let $\vec F=(z,\,x,\,y)$ and let $S$ be the part of the paraboloid $z=1-x^2-y^2$
above the $xy$-plane, oriented upward, with boundary the unit circle. Use Stokes'
theorem to compute $\oint_{\partial S}\vec F\cdot\dd\vec r$ by integrating
$\operatorname{curl}\vec F\cdot\vec n$ over the flat unit disk instead (legal because
only the rim matters).
\end{exercise}

\begin{exercise}
\label{ch35:exr-hole}
Consider $\vec F=\big(\frac{-y}{x^2+y^2},\,\frac{x}{x^2+y^2}\big)$. (a) Verify that
its scalar curl is $0$ wherever it is defined. (b) Compute its circulation around
the unit circle and get $2\pi$. (c) Explain why parts (a) and (b) do not contradict
Green's theorem. (What goes wrong at the origin?)
\end{exercise}

\begin{exercise}
\label{ch35:exr-unify}
In your own words, state the single sentence that contains the FTC, the gradient
theorem, Green's, Stokes', and the divergence theorem at once. Identify, for each of
the five, what plays the role of ``the inside,'' ``the derivative taken inside,''
and ``the boundary.''
\end{exercise}

\begin{exercise}[From field to trajectory]
\label{ch35:exr-ode-bridge}
For the swirl field $\vec F(x,y) = (-y, x)$, write down the corresponding system of
differential equations for a moving particle: $\dv{x}{t} = P(x,y)$ and $\dv{y}{t} = Q(x,y)$.
Show by direct differentiation that $\vec r(t) = \big(R\cos(t-t_0),\, R\sin(t-t_0)\big)$
solves this system for any choice of radius $R$ and initial phase $t_0$. How does this
family of solution curves relate to the concentric circles of Figure~\ref{ch35:fig-swirl}?
\end{exercise}
